\documentclass{aastex631}

\usepackage{amsmath,amssymb}
\usepackage{natbib}

\newcommand{\dd}{\mathrm{d}}
\newcommand{\pp}{\partial}
\newcommand{\kbbar}{k_{\rm B}/\mathrm{amu}}
\newcommand{\Hp}{H_P}
\newcommand{\Hpress}{H_P}
\newcommand{\Cp}{C_P}
\newcommand{\Cv}{C_V}
\newcommand{\Tint}{T_{\rm int}}
\newcommand{\Teff}{T_{\rm eff}}
\newcommand{\Tsurf}{T_{\rm surf}}
\newcommand{\Tmelt}{T_{\rm melt}}
\newcommand{\Ttrial}{T_{\rm trial}}
\newcommand{\Yeq}{Y_{\rm misc}}
\newcommand{\Zeq}{Z_{\rm misc}}
\newcommand{\softplus}{\mathrm{softplus}}

\shorttitle{ORCHARD Transport Methods}
\shortauthors{Tejada Arevalo}

\begin{document}

\title{Thermal and Compositional Transport in the ORCHARD Planetary Evolution Code}

\author{Roberto Tejada Arevalo}

\begin{abstract}
We describe the thermal and compositional transport methodology implemented
in the ORCHARD planetary evolution code.  The solver simultaneously updates
the specific entropy $S$, helium mass fraction $Y$, and heavy-element mass
fraction $Z$ on a one-dimensional Lagrangian mass grid using a fully
implicit, global Newton--Raphson scheme.  Heat transport includes
mixing-length convection, radiative diffusion, conduction, semiconvection,
and latent heat release.  Composition transport couples diffusion and
advective sedimentation driven by hydrogen--helium and hydrogen--water
immiscibility.  A smooth melt-fraction formalism bridges the inviscid
liquid and high-viscosity solid regimes in the rocky mantle, yielding a
self-consistent treatment of viscous convection from fully molten to fully
solidified states.  The coupled $3N\times 3N$ system is solved with a
trust-region--damped Newton iteration using sparse linear algebra.
\end{abstract}

\keywords{planets and satellites: interiors --- planets and satellites:
physical evolution --- methods: numerical}

% ======================================================================
\section{Introduction}\label{sec:intro}
% ======================================================================

ORCHARD evolves planetary interiors by coupling hydrostatic structure
equations with thermal and compositional transport.  The transport module
takes as input the current mechanical state of the planet---temperature
$T(m)$, pressure $P(m)$, density $\rho(m)$, and radii $r(m)$---and
returns updated profiles of entropy $S$, helium fraction $Y$, and
heavy-element fraction $Z$.  These profiles are then fed back to the
hydrostatic solver for self-consistent iteration
\citep{Henyey1964,Sur2024}.

The planet is divided into three structural regions indexed by mass
coordinate $m$:
\begin{enumerate}
  \item \textbf{Envelope} (cells $0$ to $k_{\rm core}-1$): hydrogen--helium
    mixture with dissolved heavy elements.
  \item \textbf{Rocky mantle} (cells $k_{\rm core}$ to $k_{\rm core,Fe}-1$):
    silicate material (e.g.\ MgSiO$_3$).
  \item \textbf{Iron core} (cells $k_{\rm core,Fe}$ to $N-1$): metallic iron.
\end{enumerate}
Special cases include coreless gas giants ($k_{\rm core}=N$) and bare-rock
planets ($k_{\rm core}=0$).

All quantities are defined on a cell-centered grid of $N$ mass shells.
Fluxes are evaluated at the $N+1$ cell boundaries (subscript~$b$) or the
$N-1$ interior boundaries (subscript~$i$).  Geometric and arithmetic means
map cell-centered values to interfaces where needed.

% ======================================================================
\section{Governing Equations}\label{sec:governing}
% ======================================================================

\subsection{Entropy Evolution}\label{sec:entropy}

The entropy of each mass shell evolves according to the first law of
thermodynamics in Lagrangian form:
\begin{equation}\label{eq:entropy_master}
  \frac{\pp S}{\pp t}
  = -\frac{1}{T}\frac{1}{\dd m}\left[
      F_{\rm conv} + F_{\rm rad} + F_{\rm semi}
    \right]
  + \frac{L}{T}\frac{\pp\chi}{\pp t}
  + \frac{Q_{\rm radio}}{T}
  - \frac{1}{T}\left[
      \frac{\pp U}{\pp Y}\bigg|_{S,\rho}\frac{\pp Y}{\pp t}
    + \frac{\pp U}{\pp Z}\bigg|_{S,\rho}\frac{\pp Z}{\pp t}
    \right],
\end{equation}
where $F_{\rm conv}$, $F_{\rm rad}$, and $F_{\rm semi}$ are the
convective, radiative/conductive, and semiconvective heat fluxes
(Section~\ref{sec:heat}); $L$ is the specific latent heat and $\chi$ the
liquid melt fraction (Section~\ref{sec:latent}); $Q_{\rm radio}$ is the
radiogenic heating rate (Section~\ref{sec:radiogenic}); and the final
bracket captures the chemical-potential coupling when composition evolves
simultaneously.

\subsection{Helium Evolution}\label{sec:helium_gov}

The helium mass fraction obeys an advection--diffusion equation:
\begin{equation}\label{eq:Y_master}
  \frac{\pp Y}{\pp t}
  = \frac{1}{\dd m}\left[
      4\pi r^2\rho\,D_Y\frac{\pp Y}{\pp r}
    \right]
  + \frac{1}{\dd m}\left[
      4\pi r^2\rho\,v_{\rm rain}\,\mathcal{A}_Y
    \right],
\end{equation}
where $D_Y$ is the effective helium diffusion coefficient
(Section~\ref{sec:comp_diff}), $v_{\rm rain}=D_Y/(\Hp\,f_{\rm thres})$ is
the sedimentation pseudo-velocity, and $\mathcal{A}_Y$ is a smooth
activation product gating the advective flux to the immiscible region
(Section~\ref{sec:helium_rain}).

\subsection{Heavy-Element Evolution}\label{sec:Z_gov}

The metallicity $Z$ evolves analogously to $Y$, but with two miscibility
curves $Z_{1,{\rm misc}}$ and $Z_{2,{\rm misc}}$ bounding the immiscible
region:
\begin{equation}\label{eq:Z_master}
  \frac{\pp Z}{\pp t}
  = \frac{1}{\dd m}\left[
      4\pi r^2\rho\,D_Z\frac{\pp Z}{\pp r}
    \right]
  + \frac{1}{\dd m}\left[
      4\pi r^2\rho\,v_{{\rm rain},Z}\,\mathcal{A}_Z
    \right].
\end{equation}
Metal rain (Section~\ref{sec:metal_rain}) follows either the
\citet{Rogers2025} rock-rain or \citet{Gupta2024} water-rain miscibility
formalism.

% ======================================================================
\section{Heat Transport}\label{sec:heat}
% ======================================================================

\subsection{Convective Flux (Mixing-Length Theory)}\label{sec:mlt}

In convectively unstable regions the heat flux follows
mixing-length theory (MLT):
\begin{equation}\label{eq:Fconv}
  F_{\rm conv}
  = -\varphi\left|\frac{\dd S_{\rm eff}}{\dd r}\right|^{n_{\rm conv}},
\end{equation}
where the effective entropy gradient includes the Ledoux composition
terms:
\begin{equation}\label{eq:bracket}
  \frac{\dd S_{\rm eff}}{\dd r}
  \equiv -\frac{\dd S}{\dd r}
  + \frac{\pp S}{\pp Y}\bigg|_{\rho,P}\frac{\dd Y}{\dd r}
  + \frac{\pp S}{\pp Z}\bigg|_{\rho,P}\frac{\dd Z}{\dd r}.
\end{equation}
A soft-hinge function clamps this bracket to non-negative values before
raising it to the power $n_{\rm conv}$.  Two smoothing options are
available: a ``classic'' cubic hinge $f(x) = x^3/(a^2+x^2)$, or a
softplus function $f(x)=\ln(1+e^{\beta x})/\beta$.

The transport prefactor $\varphi$ for the \textbf{envelope} is
\begin{equation}\label{eq:phi_env}
  \varphi_{\rm env}
  = \sqrt{\delta}\;4\pi r^2\,\rho\,T\,
    \sqrt{\frac{g\,\alpha_{\rm mlt}^4}{32\,\Cp}},
\end{equation}
where $\delta=-(\pp\ln\rho/\pp\ln T)_P$ is the thermal expansion
parameter from the equation of state, $g=Gm/r^2$ is the local
gravitational acceleration, $\alpha_{\rm mlt}$ is the mixing-length
parameter (denoted $\Hp$ in the code, with units of cm), and $\Cp$ is the
specific heat at constant pressure.  The convective exponent in the
envelope is $n_{\rm conv}=3/2$.

\subsection{Radiative and Conductive Flux}\label{sec:rad}

The diffusive (radiative plus conductive) heat flux is
\begin{equation}\label{eq:Frad}
  F_{\rm rad} = -\Lambda\,\frac{\dd T}{\dd r},
\end{equation}
where $\Lambda = \Lambda_{\rm cond} + \Lambda_{\rm rad}$ is the total
thermal diffusion coefficient.

\paragraph{Radiative opacity.}
In the envelope, the radiative contribution is
\begin{equation}\label{eq:Lambda_rad}
  \Lambda_{\rm rad}
  = \frac{4aT^3}{3\kappa_R\rho}\,\Theta(t),
  \qquad
  \Theta(t) = \frac{2}{\pi}\arctan\!\left(\frac{t}{t_{\rm rad}}\right),
\end{equation}
where $\kappa_R(\rho,T,Z)$ is the Rosseland mean opacity interpolated from
pre-tabulated grids at metallicities $[5,10,30,50,100]\times Z_\odot$, $a$
is the radiation density constant, and $\Theta(t)$ is a temporal ramp
(with $t_{\rm rad}=0.5$~Gyr) that prevents radiative transport from
dominating at very early times.

\paragraph{Envelope conduction.}
The conductive contribution in the envelope uses the \citet{French2019}
water conductivity as a proxy for the heavy-element component:
\begin{equation}\label{eq:Lambda_cond_env}
  \Lambda_{\rm cond}
  = \lambda_{\rm fac}\left[
      (1-Z)\,\Lambda_{\rm H/He}(P)
    + Z\,\Lambda_{\rm H_2O}(\rho,T)
    \right],
\end{equation}
where $\Lambda_{\rm H/He}(P)$ is an anchored baseline derived from
\citet{French2019} and \citet{Becker2018} adiabatic conductivity tables,
corrected for the heavy-element content at the reference adiabat:
\begin{equation}
  \Lambda_{\rm H/He}
  = \frac{\Lambda_{\rm ref}(P) - Z_{\rm ref}(P)\,\Lambda_{\rm H_2O}}
         {\max(1 - Z_{\rm ref},\,\epsilon)}.
\end{equation}
The factor $\lambda_{\rm fac}$ is a user-configurable scaling.

\paragraph{Mantle conduction.}
In the rocky mantle the thermal conductivity combines lattice and
electronic contributions:
\begin{equation}\label{eq:k_mantle}
  k_{\rm mantle} = k_{\rm lat} + k_{\rm elec}.
\end{equation}
The lattice component follows \citet{PengDeng2025} (or
\citealt{DengStixrude2021}):
\begin{equation}\label{eq:k_lat}
  k_{\rm lat}
  = k_0\left(\frac{V_0}{V}\right)^g\left(\frac{T_0}{T}\right)^a,
\end{equation}
with reference molar volume $V_0$, conductivity $k_0$, and exponents $g$
and $a$.

The electronic component uses the Wiedemann--Franz law:
\begin{equation}\label{eq:k_elec}
  k_{\rm elec} = \mathcal{L}\,\sigma(P,T)\,T,
\end{equation}
where $\mathcal{L}=2.7\times10^{-8}$~W\,$\Omega$\,K$^{-2}$ is the Lorenz
number and $\sigma$ is the electrical conductivity from
\citet{Holmstrom2018}:
\begin{equation}\label{eq:sigma}
  \sigma = \sigma_0\exp\!\left[
    -\frac{\Delta E + P\Delta V}{RT}
  \right],
\end{equation}
with $\sigma_0=1.99\times10^5$~S\,m$^{-1}$,
$\Delta E=75.94$~kJ\,mol$^{-1}$, and
$\Delta V=-0.061$~cm$^3$\,mol$^{-1}$.

\paragraph{Iron core.}
The iron core uses a constant thermal conductivity (user-specified,
typically 40--100~W\,m$^{-1}$\,K$^{-1}$).

\subsection{Semiconvective Flux}\label{sec:semi}

Semiconvection operates in regions that are Schwarzschild-unstable but
Ledoux-stable.  The double-diffusive criterion is parameterized by the
inverse density ratio:
\begin{equation}\label{eq:R0inv}
  R_0^{-1}
  = \frac{\nabla_S - \nabla_L}{\nabla_S},
\end{equation}
where $\nabla_S$ and $\nabla_L$ are the Schwarzschild and Ledoux entropy
gradients, respectively.  Semiconvection is active where
$1 < R_0^{-1} < R_{\rm crit}^{-1}$, with
\begin{equation}
  R_{\rm crit}^{-1} = \frac{\Pr + 1}{\Pr + \tau},
\end{equation}
where $\Pr$ is the turbulent Prandtl number and $\tau$ the
diffusivity ratio.

\subsubsection{Continuous Model}

The default (continuous) semiconvection model maps the double-diffusive
window onto a stability parameter:
\begin{equation}
  \chi_{\rm sc}
  = \frac{R_0^{-1} - 1}{R_{\rm crit}^{-1} - 1},
  \qquad 0\le\chi_{\rm sc}\le 1,
\end{equation}
and defines a smooth ramp
\begin{equation}
  f_{\rm sc} = (1 - \chi_{\rm sc})^{n_{\rm smooth}},
\end{equation}
which modulates the thermal Nusselt number:
\begin{equation}
  \mathrm{Nu}_T = 1 + (\mathrm{Nu}_{T,\max} - 1)\,f_{\rm sc}.
\end{equation}
The semiconvective heat flux is then
\begin{equation}\label{eq:Fsemi}
  F_{\rm semi}
  = -4\pi r^2\,\Lambda\,\frac{(\mathrm{Nu}_T - 1)\,T}{\Cp}
    \,\nabla_S,
\end{equation}
and the compositional diffusivity is enhanced by a matching Nusselt
number $\mathrm{Nu}_X$:
\begin{equation}
  D_{\rm semi} = D_{\rm micro}\,\mathrm{Nu}_X.
\end{equation}
Both $\mathrm{Nu}_T$ and $\mathrm{Nu}_X$ are time-lagged via exponential
relaxation to avoid Newton-iteration flicker:
\begin{equation}
  \mathrm{Nu}(t+\Delta t)
  = \mathrm{Nu}_{\rm old}
  + \theta_\nu\,(\mathrm{Nu}_{\rm tgt} - \mathrm{Nu}_{\rm old}),
  \quad
  \theta_\nu = 1 - e^{-\Delta t/\tau_D}.
\end{equation}

\subsubsection{Bulk (Semi-Shoot) Model}

An alternative ``semi-shoot'' model applies binary semiconvection in
zones satisfying the $R_0^{-1}$ criterion and extends it with an
exponential overshoot taper around the envelope--mantle boundary:
\begin{equation}
  w_{\rm ov}(r)
  = \exp\!\left(-\frac{|r - r_{\rm EMB}|}{f_{\Hp}\,\Hp}\right),
\end{equation}
where $r_{\rm EMB}$ is the radius of the envelope--mantle interface and
$f_{\Hp}$ controls the overshoot scale in units of the local pressure
scale height.

% ======================================================================
\section{Viscous Mantle Convection}\label{sec:viscous}
% ======================================================================

The rocky mantle transitions from fully liquid (low viscosity,
$\eta_{\rm liq}\sim100$~Pa\,s) to fully solid (high viscosity,
$\eta_{\rm solid}(P,T)$ from dislocation creep) as it cools.  ORCHARD
models this transition using a smooth melt-fraction parameter $\chi$
(Section~\ref{sec:latent}) that blends between two convective regimes.

\subsection{Solid-State Viscosity}

The dynamic viscosity of solid silicate follows a dislocation-creep law
\citep{Ranalli2001}:
\begin{equation}\label{eq:eta_solid}
  \eta_{\rm solid}
  = \frac{1}{2}\,B^{-1/n}\,\dot\varepsilon^{(1-n)/n}
    \exp\!\left(\frac{E^* + PV^*}{nRT}\right),
\end{equation}
where $B$, $n$, $E^*$, and $V^*$ are material-dependent parameters
(perovskite values: $B=7.4\times10^{-17}$, $n=3.5$,
$E^*=5.0\times10^5$~J\,mol$^{-1}$, $V^*=10^{-5}$~m$^3$\,mol$^{-1}$)
and $\dot\varepsilon=10^{-15}$~s$^{-1}$ is the assumed strain rate.

The effective kinematic viscosity at each boundary is blended in
log-space:
\begin{equation}\label{eq:nu_blend}
  \log_{10}\nu_{\rm eff}
  = \chi\,\log_{10}\nu_{\rm liq}
  + (1-\chi)\,\log_{10}\nu_{\rm solid},
\end{equation}
where $\nu = \eta/\rho$.

\subsection{Eddy Diffusivity Branches}

Four eddy-diffusivity branches capture the physics:

\paragraph{Inviscid, thermal (Eq.~34a):}
\begin{equation}\label{eq:kappa34a}
  \kappa_{34a}
  = \sqrt{\frac{\alpha\,T\,g\,\Hp^4}{32\,\Cp}}.
\end{equation}

\paragraph{Viscosity-limited, thermal (Eq.~34b):}
\begin{equation}\label{eq:kappa34b}
  \kappa_{34b}
  = \frac{\alpha\,T\,g\,\Hp^4}{32\,\Cp\,\nu_{\rm eff}}.
\end{equation}

\paragraph{Inviscid, melt-curve (Eq.~35a):}
\begin{equation}\label{eq:kappa35a}
  \kappa_{35a}
  = \sqrt{\frac{\alpha_\chi\,\Tmelt\,g\,\Hp^4}{32\,L}},
\end{equation}
where $\alpha_\chi = |(\pp\rho/\pp\chi)_{P,T}|/\rho$ is the
melt-contraction parameter and $L$ is the specific latent heat.

\paragraph{Viscosity-limited, melt-curve (Eq.~35b):}
\begin{equation}\label{eq:kappa35b}
  \kappa_{35b}
  = \frac{\alpha_\chi\,\Tmelt\,g\,\Hp^4}{32\,L\,\nu_{\rm eff}}.
\end{equation}

\subsection{Log-Space Blending}

The effective eddy diffusivity interpolates between the inviscid and
viscous branches using the melt fraction as a weight.  First, the
composite inviscid and viscous diffusivities are formed:
\begin{align}
  \kappa_{\rm inv}
  &= \chi\,\kappa_{34a} + (1-\chi)\,\kappa_{35a},
  \label{eq:kappa_inv}\\
  \kappa_{\rm vis}
  &= (1-\chi)\,\kappa_{34b} + \chi\,\kappa_{35b}.
  \label{eq:kappa_vis}
\end{align}
The dynamic range of the ratio
$\log_{10}(\kappa_{\rm inv}/\kappa_{\rm vis})$ is capped at
$\Delta_{\max}=6$ to prevent catastrophic sensitivity near the melt front.
The final diffusivity is
\begin{equation}\label{eq:kappa_eff}
  \log_{10}\kappa_{\rm eff}
  = \chi\,\log_{10}\kappa_{\rm inv}
  + (1-\chi)\,\log_{10}\kappa_{\rm vis}.
\end{equation}

The mantle transport prefactor is then
\begin{equation}\label{eq:phi_mantle}
  \varphi_{\rm mantle}
  = 4\pi r^2\,\rho\,T\,\kappa_{\rm eff},
\end{equation}
and the convective exponent smoothly transitions with melt fraction:
\begin{equation}\label{eq:nconv_mantle}
  n_{\rm conv}(\chi)
  = \tfrac{3}{2} + \tfrac{1}{2}(1-\chi),
\end{equation}
i.e.\ $n=3/2$ when fully liquid and $n=2$ when fully solid.

% ======================================================================
\section{Latent Heat}\label{sec:latent}
% ======================================================================

\subsection{Melt Fraction}

The liquid melt fraction $\chi(T,P)$ is modeled as a smooth step:
\begin{equation}\label{eq:chi}
  \chi(T,P)
  = \frac{1}{2}\left[1 + \tanh\!\left(
      \frac{T - \Tmelt(P)}{\Delta T_{\rm lat}}
    \right)\right],
\end{equation}
where $\Tmelt(P)$ is the pressure-dependent melting curve (from the
mantle or core equation of state) and $\Delta T_{\rm lat}$ controls the
width of the transition (default 300~K).  Its derivative is
\begin{equation}\label{eq:dchidT}
  \frac{\pp\chi}{\pp T}
  = \frac{1}{2\,\Delta T_{\rm lat}}
    \,\mathrm{sech}^2\!\left(
      \frac{T - \Tmelt}{\Delta T_{\rm lat}}
    \right).
\end{equation}

\subsection{Relaxation}\label{sec:chi_relax}

The equilibrium melt fraction $\chi_{\rm tgt}(T,P)$ from
Equation~(\ref{eq:chi}) can change abruptly between timesteps when the
temperature crosses the melting curve, particularly in narrow transition
zones where $\Delta T_{\rm lat}$ is small relative to the temperature
change per step.  To prevent such discontinuous jumps from
destabilizing the Newton solver, a first-order relaxation filter is
applied:
\begin{equation}\label{eq:chi_relax}
  \chi^{n+1}
  = \chi^n + \theta_\chi\,\bigl(\chi_{\rm tgt}(T_{\rm trial}) - \chi^n\bigr),
\end{equation}
where $\chi^n$ is the melt fraction accepted at the previous timestep,
$\chi_{\rm tgt}(T_{\rm trial})$ is the equilibrium target evaluated at
the \emph{trial temperature} (Section~\ref{sec:trial_temp}) of the
current Newton iterate, and
\begin{equation}\label{eq:theta_chi}
  \theta_\chi = 1 - e^{-\Delta t/\tau_\chi}
\end{equation}
is a dimensionless blending weight controlled by the user-specified
relaxation timescale $\tau_\chi$.

The behavior of the filter depends on the ratio $\Delta t/\tau_\chi$:
\begin{itemize}
  \item \textbf{Small timestep} ($\Delta t \ll \tau_\chi$):
    $\theta_\chi \approx \Delta t/\tau_\chi$, so $\chi$ evolves at a
    rate limited by $\tau_\chi$ regardless of how large the
    equilibrium jump is:
    \begin{equation}\label{eq:dchidt_limit}
      \frac{\Delta\chi}{\Delta t}
      \approx \frac{\chi_{\rm tgt} - \chi^n}{\tau_\chi}\,.
    \end{equation}
  \item \textbf{Large timestep} ($\Delta t \gg \tau_\chi$):
    $\theta_\chi \to 1$, and $\chi^{n+1} \to \chi_{\rm tgt}$---the
    full equilibrium value is recovered.
  \item \textbf{No relaxation} ($\tau_\chi = 0$):
    $\theta_\chi = 1$ identically, and
    $\chi^{n+1} = \chi_{\rm tgt}(T_{\rm trial})$.
\end{itemize}

Two separate relaxation timescales are maintained: one for the
latent-heat source term (which directly enters the entropy residual)
and one for the transport-geometry updates (which control the
viscous/inviscid blending in mantle convection,
Section~\ref{sec:mantle}).  This allows the viscous transition to
track the equilibrium melt state more or less tightly than the
latent-heat release.

\subsection{Nature of $\Delta\chi$}\label{sec:dchi_nature}

The term $\pp\chi/\pp t$ in the continuous entropy equation
(Equation~\ref{eq:entropy_master}) is discretized as a finite
difference $\Delta\chi / \Delta t$, but it is \emph{not} a simple
forward-Euler difference of the equilibrium melt fraction.  Instead,
$\Delta\chi$ is the change in the \textbf{relaxation-filtered} melt
fraction evaluated at the \textbf{trial temperature}:
\begin{equation}\label{eq:dchi_def}
  \Delta\chi_k
  \;\equiv\; \chi_k^{n+1} - \chi_k^n
  \;=\; \theta_\chi \bigl(\chi_{\rm tgt}(T_{{\rm trial},k}) - \chi_k^n\bigr).
\end{equation}
Two features make this implicit:
\begin{enumerate}
  \item \textbf{Trial-temperature dependence.}  The target
    $\chi_{\rm tgt}$ is evaluated at $T_{\rm trial}$, which is a
    linearized function of the Newton unknowns $(S', Y, Z)$ via
    \begin{equation}
      T'_{\rm trial} = T' + \frac{T'}{\Cv}\,S_0\,(S' - S'_{\rm old})
        + \frac{\pp T'}{\pp Y}\,(Y - Y_{\rm old})
        + \frac{\pp T'}{\pp Z}\,(Z - Z_{\rm old}).
    \end{equation}
    Therefore $\Delta\chi$ changes as the Newton solver updates $S$,
    $Y$, and $Z$, and the Jacobian captures this coupling.
  \item \textbf{Relaxation damping.}  The prefactor $\theta_\chi$
    limits how much $\chi$ can change per timestep, preventing the
    latent heat from releasing or absorbing energy faster than the
    relaxation timescale permits.
\end{enumerate}

\subsection{Residual Contribution}\label{sec:lat_residual}

The latent heat enters the entropy residual
(Equation~\ref{eq:BS}) as
\begin{equation}\label{eq:Rlat}
  \mathcal{R}_{{\rm lat},k}
  = \frac{L_k}{T_0\,S_0}\,\frac{\Delta\chi_k}{T'_k}\,,
\end{equation}
where $L_k$ is the specific latent heat of the relevant layer
(mantle or iron core), $T'_k = T_k/T_0$ is the dimensionless
temperature, and $\Delta\chi_k$ is given by
Equation~(\ref{eq:dchi_def}).  When $\Delta\chi > 0$ (melting), the
residual drives the entropy \emph{downward} (the material absorbs
latent heat from its thermal energy); when $\Delta\chi < 0$
(freezing), entropy is released.

\subsection{Jacobian Contributions}\label{sec:lat_jacobian}

Because $\Delta\chi_k$ depends on $T_{\rm trial}$, and therefore on
the Newton unknowns, the Jacobian of the latent residual is nonzero.
Define the common chain-rule factor
\begin{equation}\label{eq:common_lat}
  C_k
  = \frac{\pp\chi_{\rm tgt}}{\pp T}\bigg|_k
    \!\cdot\,\frac{T_0}{T'_k}
  \;-\; \frac{\Delta\chi_k}{(T'_k)^2}\,,
\end{equation}
where the first term differentiates the numerator $\Delta\chi$ with
respect to $T$ (via the tanh derivative,
Equation~\ref{eq:dchidT}, multiplied by~$\theta_\chi$), and the
second term differentiates the $1/T'$ denominator.
The diagonal Jacobian contributions in the entropy row are then
\begin{align}
  J_{{\rm lat},S,k}
    &= \frac{L_k}{T_0\,S_0}
       \,\frac{\pp T'}{\pp S'}\bigg|_k \! C_k\,,
    &\frac{\pp T'}{\pp S'}
    &= \frac{T'}{\Cv}\,S_0\,,
  \label{eq:JlatS}\\[4pt]
  J_{{\rm lat},Y,k}
    &= \frac{L_k}{T_0\,S_0}
       \,\frac{\pp T'}{\pp Y}\bigg|_k \! C_k\,,
  \label{eq:JlatY}\\[4pt]
  J_{{\rm lat},Z,k}
    &= \frac{L_k}{T_0\,S_0}
       \,\frac{\pp T'}{\pp Z}\bigg|_k \! C_k\,.
  \label{eq:JlatZ}
\end{align}
These are added to the diagonal entries of the Jacobian:
$J_{{\rm lat},S}$ enters $J_{3k,3k}$,
$J_{{\rm lat},Y}$ enters $J_{3k,3k+1}$, and
$J_{{\rm lat},Z}$ enters $J_{3k,3k+2}$.
No off-diagonal (neighbor-cell) Jacobian entries arise because $\chi$
depends only on the local cell temperature and pressure.

The key point is that the $\pp\chi/\pp T$ derivative appearing in
Equations~(\ref{eq:JlatS})--(\ref{eq:JlatZ}) is the derivative of
the \emph{relaxed} melt fraction, not the raw equilibrium value:
\begin{equation}\label{eq:dchidT_relaxed}
  \frac{\pp\chi^{n+1}}{\pp T}
  = \theta_\chi\,\frac{\pp\chi_{\rm tgt}}{\pp T}\,.
\end{equation}
When $\theta_\chi < 1$, the effective sensitivity is reduced, keeping
the Jacobian consistent with the damped residual and preventing the
Newton solver from overcorrecting near the melt front.

\subsection{Latent-Heat Parameters}

The mantle latent heat is $L_{\rm mantle}=3.3\times10^9\,f_L$~erg\,g$^{-1}$
(where $f_L$ is a user-configurable scale factor), and the core value is
provided by the iron equation of state.

% ======================================================================
\section{Radiogenic Heating}\label{sec:radiogenic}
% ======================================================================

Radioactive decay in the rocky mantle contributes a volumetric heat
source to the entropy equation.  The per-cell radiogenic power is
\begin{equation}\label{eq:Qradio}
  Q_{{\rm radio},k}(t) = H_{\rm spec}(t)\,\delta m_k,
\end{equation}
where $\delta m_k$ is the cell mass and $H_{\rm spec}(t)$
(erg\,s$^{-1}$\,g$^{-1}$) is the time-dependent specific heating
rate.  Radiogenic heating is applied only to mantle cells
($k_{\rm core}\le k < k_{\rm core,Fe}$); envelope and iron-core cells
receive no radiogenic contribution.

\subsection{Isotope Decay Chains}

The specific heating rate is the sum of contributions from four
long-lived isotopes:
\begin{equation}\label{eq:Hspec}
  H_{\rm spec}(t)
  = \sum_{i\in\mathcal{I}} w_i\,q_{0,i}\;
    2^{(t_{\rm now} - t)/\tau_i},
\end{equation}
where $\mathcal{I}=\{{}^{40}\mathrm{K},\;{}^{232}\mathrm{Th},\;
{}^{235}\mathrm{U},\;{}^{238}\mathrm{U}\}$, $q_{0,i}$ are
present-day Earth heat-production rates per unit mass from
\citet{McDonough1995}, $\tau_i$ are the respective half-lives, and
$t_{\rm now}=4.567$\,Gyr is the present Solar System age.  The
optional weight factors $w_i$ (default unity) allow the user to scale
the radiogenic inventory relative to Earth.  Because the exponential
runs backwards from today, the heating rate at early times is
significantly larger: the short-lived isotopes $^{235}$U
($\tau=0.704$\,Gyr) and $^{40}$K ($\tau=1.25$\,Gyr) dominate during
the first few billion years.  Table~\ref{tab:isotopes} lists the
adopted constants.

\begin{deluxetable}{lcc}
\tablecaption{Parameters of Radiogenic Elements \citep{McDonough1995}.\label{tab:isotopes}}
\tablecolumns{3}
\tablehead{
  \colhead{Element} &
  \colhead{$q_0$ (W\,kg$^{-1}$)} &
  \colhead{$\tau$ (Gyr)}
}
\startdata
$^{40}$K    & $8.69\times10^{-13}$ & 1.25  \\
$^{232}$Th  & $2.24\times10^{-12}$ & 14    \\
$^{235}$U   & $8.48\times10^{-14}$ & 0.704 \\
$^{238}$U   & $1.97\times10^{-12}$ & 4.47  \\
\enddata
\end{deluxetable}

\subsection{Entry into the Entropy Residual}

The radiogenic term enters the entropy residual
(Equation~\ref{eq:BS}) as
\begin{equation}\label{eq:Bsrad}
  B_{3k} \mathrel{-}= \frac{\Delta t}{\dd m'_k\,T'_k}\,Q'_k,
\end{equation}
where $Q'_k = Q_{{\rm radio},k}/(M_0\,T_0\,S_0)$ is the
non-dimensionalized per-cell power.  The minus sign ensures that
positive heating drives entropy upward under the Newton update
$\mathbf{z}_{\rm new}=\mathbf{z}_{\rm old}-\mathbf{J}^{-1}\mathbf{B}$.

\subsection{Jacobian Contribution}

Although $H_{\rm spec}$ depends only on time (not on the solution
variables $S$, $Y$, $Z$), the residual term
$-\Delta t\,Q'_k/(\dd m'_k\,T'_k)$ depends implicitly on entropy
through the $1/T'_k$ factor.  Differentiating with respect to $S'_k$
gives
\begin{equation}\label{eq:Jrad}
  \frac{\pp}{\pp S'_k}\!\left[
    -\frac{\Delta t}{\dd m'_k}\,\frac{Q'_k}{T'_k}
  \right]
  = \frac{\Delta t}{\dd m'_k}\,
    \frac{Q'_k}{(T'_k)^{2}}\,
    \frac{\pp T'}{\pp S'}\bigg|_k,
\end{equation}
where the EOS chain rule provides
\begin{equation}\label{eq:dTdS_radio}
  \frac{\pp T'}{\pp S'}\bigg|_k
  = \frac{T'_k}{\Cv[,k]}\,S_0.
\end{equation}
This is added to the diagonal Jacobian entry $J_{3k,\,3k}$.
The contribution is positive, reflecting the stabilizing feedback
that a hotter cell (higher~$S$) reduces the $Q/T$ source and thereby
slows further entropy production.  No off-diagonal Jacobian entries
arise from the radiogenic term because $H_{\rm spec}$ is independent
of composition.

% ======================================================================
\section{Composition Transport}\label{sec:comp}
% ======================================================================

\subsection{Diffusion Coefficients}\label{sec:comp_diff}

The effective diffusion coefficient for helium is the sum of a
microscopic and a convective (MLT) contribution:
\begin{equation}\label{eq:DY}
  D_Y = D_{\rm micro,Y} + D_{\rm MLT},
\end{equation}
where the MLT diffusivity is
\begin{equation}\label{eq:DMLT}
  D_{\rm MLT}
  = \frac{v_{\rm MLT}\,\Hpress}{3},
  \qquad
  v_{\rm MLT}
  = \sqrt{\frac{g\,\Hp^2}{8\,\Cp}\,\frac{\dd S_{\rm eff}}{\dd r}}.
\end{equation}
The metal diffusion coefficient $D_Z$ has the same structure but with
separate microscopic values $D_{\rm micro,Z}$.  An elevated interior
value $D_{\rm micro,Z,in}$ is used for enclosed mass fractions
$m/M < f_{\rm inner}$ (default 0.7) to prevent unrealistic deep
stratification.

\subsubsection{Smooth Ledoux Transition}

In zones that are Schwarzschild-unstable but Ledoux-stable, the
composition diffusivity is smoothly ramped down rather than cut off:
\begin{equation}
  D_{\rm trans}
  = D_{\rm Schwarz}\,(1 - \chi_{\rm stab})^{n_{\rm smooth}},
\end{equation}
where $\chi_{\rm stab}=|\nabla_L|/(\nabla_S - \nabla_L)$ measures
the Ledoux stabilization strength and $D_{\rm Schwarz}$ is the
Schwarzschild-only MLT diffusivity.

\subsubsection{Time Relaxation}

To prevent large timestep-to-timestep jumps, $D_Y$ and $D_Z$ are relaxed
using a closed-form M\"obius update:
\begin{equation}\label{eq:D_relax}
  w = \frac{D - D_A}{D - D_B},
  \qquad
  w(t+\Delta t) = w(t)\,e^{-\Delta t/(\tau_D\,|D_A D_B|)},
  \qquad
  D = \frac{D_B\,w - D_A}{w - 1},
\end{equation}
where $D_A$ and $D_B$ are the instantaneous and microscopic targets, and
$\tau_D$ is the diffusion relaxation timescale.  A stable exponential
blend serves as a numerical fallback when the M\"obius transform becomes
ill-conditioned.

\subsection{Helium Rain}\label{sec:helium_rain}

Helium sedimentation is driven by the hydrogen--helium miscibility curve
$\Yeq(P,T)$, evaluated from \citet{Lorenzen2011} or
\citet{Schoettler2018} tables.

\subsubsection{Activation Product}

The advective helium-rain flux is gated by a smooth activation product
that is nonzero only in the immiscible region:
\begin{equation}\label{eq:Yflux_new}
  \mathcal{A}_Y
  = \softplus_\beta\!\bigl(Y - \Yeq(1-Z) - \delta_{\rm imp}\bigr)
  \;\cdot\;
  \softplus_\beta\!\bigl(Y_{\rm thres} - Y\bigr),
\end{equation}
where $\softplus_\beta(x) = \ln(1+e^{\beta x})/\beta$ is a smooth ReLU
approximation, $Y_{\rm thres}$ is an inner threshold preventing
sedimentation past the core, and $\delta_{\rm imp}$ captures the implicit
temperature dependence of the miscibility curve:
\begin{equation}
  \delta_{\rm imp}
  = \left(\frac{\pp\Yeq}{\pp T}
  + \frac{\pp\Yeq}{\pp P}\frac{\pp P}{\pp T}\right)
  \frac{T}{\Cp}\,(S - S_{\rm old})\,S_0\,(1-Z).
\end{equation}

The sign of the rain flux uses a smooth-sign surrogate:
\begin{equation}
  \widetilde{\mathrm{sgn}}(x)
  = \frac{x}{\sqrt{x^2 + \epsilon^2}},
\end{equation}
with $\epsilon\sim10^{-12}$.

\subsubsection{Discretized Flux Coefficients}\label{sec:disc_comp_flux}

At each interior boundary $i$ (lying between cells $k$ and $k+1$), the
composition fluxes are constructed from a diffusive and an advective
coefficient.  Let $r_i$, $\rho_i$ denote the geometric-mean radius and
density at boundary~$i$, and $\Delta r_i = r_{k+1} - r_k$ the radial
cell spacing.  The common geometric prefactor is
$\mathcal{G}_i = 4\pi r_i^2\,\rho_i$.

The \textbf{diffusive} flux coefficient for helium is
\begin{equation}\label{eq:sigma1}
  \sigma_{1,i} = \frac{\mathcal{G}_i\,D_{Y,i}}{\Delta r_i},
\end{equation}
where $D_{Y,i}$ is the helium diffusion coefficient interpolated to
boundary~$i$.  The \textbf{advective} (sedimentation) flux coefficient is
\begin{equation}\label{eq:sigma2}
  \sigma_{2,i}
  = \frac{\mathcal{G}_i\,D_{2,i}}{\Hp\,f_{\rm thres}},
\end{equation}
where $D_{2,i}$ is the advective diffusion coefficient (equal to the
MLT diffusivity in the rain zone) and $f_{\rm thres}$ is a
miscibility-threshold parameter controlling the rain velocity.

The total helium flux at boundary~$i$ is therefore
\begin{equation}\label{eq:Yflux_total}
  \mathcal{F}_{Y,i}
  = \sigma_{1,i}\,(Y_{k+1}-Y_k)
  + \sigma_{2,i}\,\widetilde{\mathrm{sgn}}(\mathcal{A}_Y)\,\mathcal{A}_{Y,i}.
\end{equation}
Analogous expressions hold for the metal fluxes with
$D_{Z,i}$, $D_{2Z,i}$, and scale height $H_Z$ replacing $\Hp\,f_{\rm thres}$.

\subsubsection{Activation-Product Derivatives}\label{sec:activ_deriv}

The Newton solver requires the Jacobian of the activation product
$\mathcal{A}_Y$ (Equation~\ref{eq:Yflux_new}) with respect to the state
variables $S'$, $Y$, and $Z$.  Define the shorthand
\begin{equation}
  \mathcal{P}_1 \equiv \softplus_\beta(\mathrm{arg}_1),
  \qquad
  \mathcal{P}_2 \equiv \softplus_\beta(\mathrm{arg}_2),
\end{equation}
where $\mathrm{arg}_1 = Y - \Yeq(1-Z) - \delta_{\rm imp}$ and
$\mathrm{arg}_2 = Y_{\rm thres} - Y$.  The derivative of the softplus
is the sigmoid function:
\begin{equation}
  \sigma_\beta(x)
  \equiv \frac{\dd}{\dd x}\softplus_\beta(x)
  = \frac{1}{1 + e^{-\beta x}}.
\end{equation}

By the product rule, the $Y$-derivative of the activation product at
boundary~$i$ is
\begin{equation}\label{eq:dAY_dY}
  \frac{\pp\mathcal{A}_Y}{\pp Y}\bigg|_i
  = \frac{1}{2}\bigl[
      \sigma_\beta(\mathrm{arg}_1)\,\mathcal{P}_2
    - \mathcal{P}_1\,\sigma_\beta(\mathrm{arg}_2)
  \bigr],
\end{equation}
where the factor $\tfrac{1}{2}$ arises from the arithmetic-mean
interpolation $Y_i = (Y_k + Y_{k+1})/2$, so
$\pp Y_i/\pp Y_k = \tfrac{1}{2}$.

The $S'$-derivative propagates through the miscibility-curve temperature
dependence:
\begin{equation}\label{eq:dAY_dS}
  \frac{\pp\mathcal{A}_Y}{\pp S'}\bigg|_i
  = -\sigma_\beta(\mathrm{arg}_1)\,\mathcal{P}_2
    \;\left(\frac{\pp\Yeq}{\pp T}
    + \frac{\pp\Yeq}{\pp P}\frac{\pp P}{\pp T}\right)_{\!i}
    \frac{T_i}{\Cp[,i]}\,\frac{S_0}{2},
\end{equation}
where the factor $(T_i/\Cp[,i])\,S_0$ converts an entropy perturbation
to a temperature perturbation via the trial-temperature linearization.
In practice, the $\pp P/\pp T$ contribution is set to zero for
convergence stability, so the miscibility-curve sensitivity enters
only through $\pp\Yeq/\pp T$.

An analogous structure holds for $\mathcal{A}_Z$, with each
$\softplus$ factor differentiated via the chain rule through
$\pp Z_{1,{\rm misc}}/\pp T$ and $\pp Z_{2,{\rm misc}}/\pp T$.

\subsubsection{Rain Models}

Two advective models are available:

\begin{itemize}
  \item \textbf{Legacy}: The diffusion coefficient is set equal to
    the MLT diffusivity in the rain zone; the activation product
    $\mathcal{A}_Y$ naturally limits the flux.
  \item \textbf{Bottom-up} \citep{Helled2025}: Identical advection
    coefficients, but convective diffusion is smoothly suppressed in
    the rain zone via a factor
    \begin{equation}
      f_{\rm supp}
      = \frac{1}{1 + \alpha_{\rm rain}\,\softplus_\beta(\Delta Y)},
    \end{equation}
    where $\Delta Y = Y - \Yeq(1-Z)$ is the supersaturation and
    $\alpha_{\rm rain}$ controls suppression strength.
\end{itemize}

\subsection{Metal Rain}\label{sec:metal_rain}

Metal rain follows the same advection--diffusion structure as helium rain,
but uses two miscibility curves $Z_{1,{\rm misc}}(P,T)$ and
$Z_{2,{\rm misc}}(P,T)$ from either \citet{Rogers2025} (rock rain) or
\citet{Gupta2024} (water rain):
\begin{align}
  \mathcal{A}_Z
  &= \softplus_\beta(Z - Z_{1,{\rm misc}} - \delta_{Z,1})
  \;\cdot\;
  \softplus_\beta(Z_{2,{\rm misc}} - Z + \delta_{Z,2}).
\end{align}
The Jacobian includes chain-rule derivatives
$\pp Z_{1,2,\rm misc}/\pp T$ and $\pp Z_{1,2,\rm misc}/\pp P$ for
implicit coupling.

\subsection{Time-Lagged Miscibility}

Both helium and metal miscibility profiles are time-lagged to smooth
the response:
\begin{equation}\label{eq:Ym_lag}
  \Yeq^{\rm lag}
  = \frac{\Yeq^{\rm old}/\Delta t + \Yeq/\tau_m}{1/\Delta t + 1/\tau_m},
\end{equation}
with an analogous expression for $Z_{1,2,\rm misc}^{\rm lag}$.

% ======================================================================
\section{Boundary Conditions}\label{sec:bc}
% ======================================================================

\subsection{Atmospheric Boundary}

The outer boundary condition couples the interior to an atmosphere model
via the intrinsic temperature $\Tint$:
\begin{equation}
  F_{\rm surf} = \sigma_{\rm SB}\,\Tint^4,
\end{equation}
where $\Tint$ is a function of surface entropy, gravity, and composition.
Multiple atmosphere tables are supported (C23/C26: \citealt{Chen2023};
F11: \citealt{Fortney2011}; B97: \citealt{Burrows1997}; F07:
\citealt{Fortney2007}), each providing derivatives
$\pp\Tint/\pp S$, $\pp\Tint/\pp Y$, $\pp\Tint/\pp Z$ for the Jacobian.

In the discretized residual, the surface radiative loss enters as a
flux at the outermost boundary ($b=0$) of cell $k=0$:
\begin{equation}\label{eq:Fsurf_disc}
  F^{(0)}_{\rm surf}
  = -\frac{A_{\rm geom}\,\sigma_{\rm SB}\,\Tint^4}{M_0\,T_0\,S_0},
\end{equation}
where $A_{\rm geom} = 4\pi R_{\rm surf}^2$ is the emitting area,
and $M_0$, $T_0$, $S_0$ are the normalization constants defined in
Section~\ref{sec:dimless}.  The corresponding Jacobian contributions
to the entropy equation of the surface cell are obtained by
differentiating Equation~(\ref{eq:Fsurf_disc}) with respect to the
state variables.  Define the prefactor
\begin{equation}\label{eq:surf_pref}
  p = \frac{\Delta t}{\dd m'_0\,T'_0}
      \;\frac{4\,A_{\rm geom}\,\sigma_{\rm SB}\,\Tint^3}
             {M_0\,T_0\,S_0}.
\end{equation}
Then the surface Jacobian entries are
\begin{align}
  J_{0,0}   &\mathrel{+}= p\;\frac{\pp\Tint}{\pp S},
  \label{eq:Jsurf_S}\\
  J_{0,1}   &\mathrel{+}= p\;\frac{\pp\Tint}{\pp Y},
  \label{eq:Jsurf_Y}\\
  J_{0,2}   &\mathrel{+}= p\;\frac{\pp\Tint}{\pp Z},
  \label{eq:Jsurf_Z}
\end{align}
where the $Y$ and $Z$ entries are included only when composition
coupling is active.

\subsection{Bare-Rock Surface}

For planets without a gaseous envelope ($k_{\rm core}=0$), the surface
cell radiates directly:
\begin{equation}
  F_{\rm int} = \sigma_{\rm SB}(\Tsurf^4 - T_{\rm eq}^4),
\end{equation}
and the Jacobian is
\begin{equation}
  \frac{\pp F_{\rm int}}{\pp S}
  = \frac{\pp F_{\rm int}}{\pp\Tsurf}
    \,\frac{T_{\rm surf}}{C_{P,{\rm surf}}}
    \,S_0.
\end{equation}
The Jacobian prefactor for the bare-rock surface is
\begin{equation}\label{eq:pref_bare}
  p_{\rm bare}
  = \frac{\Delta t}{\dd m'_0\,T'_0}
    \;\frac{A_{\rm geom}}{M_0\,T_0\,S_0},
\end{equation}
and the entropy Jacobian entry is
$J_{0,0} \mathrel{+}= p_{\rm bare}\,(\pp F_{\rm int}/\pp S)$,
where
$\pp F_{\rm int}/\pp S
 = (\pp F_{\rm int}/\pp\Tsurf)\,(T_{\rm surf}/\Cp[,\rm surf])\,S_0$.
Semi-gray variants (``gray'' and ``ideal gray'') are also available,
using an optical-depth correction $\tau_{\rm match}$ at the matching
pressure, so that the intrinsic flux becomes
$F_{\rm int} = \sigma_{\rm SB}\,\Tsurf^4/[\tfrac{3}{4}(\tau_{\rm match}+\tfrac{2}{3})]$.

\subsection{Bottom Boundary}\label{sec:bc_bottom}

At the innermost boundary ($b=N$), an insulating (zero-flux) condition
is applied.  This is enforced implicitly: all flux arrays are initialized
to zero and populated only at interior indices $1$ through $N-1$, so
$F_N = 0$ identically.  No Jacobian entries are generated for the
innermost boundary.

% ======================================================================
\section{Numerical Method}\label{sec:numerics}
% ======================================================================

\subsection{Dimensionless Variables}\label{sec:dimless}

All quantities entering the Newton solver are cast in dimensionless
form to improve the conditioning of the $3N\times 3N$ linear system.
Four reference scales are introduced:
\begin{align}
  S_0 &= \kbbar
       &&\text{(entropy scale)},
  \label{eq:S0}\\
  T_0 &= 1000\;\mathrm{K}
       &&\text{(temperature scale)},
  \label{eq:T0}\\
  M_0 &= M_{\rm planet}
       &&\text{(mass scale)},
  \label{eq:M0}\\
  R_0 &= R_{\rm planet}
       &&\text{(radius scale)}.
  \label{eq:R0}
\end{align}
$M_0$ and $R_0$ are set to the planet's total mass and initial radius
(e.g.\ $M_{\rm Jup}$ and $R_{\rm Jup}$ for a Jupiter model).
The dimensionless state variables are
\begin{equation}\label{eq:dimless_vars}
  S' = \frac{S\,S_0}{\tilde{S}_0},
  \qquad
  T' = \frac{T}{T_0},
  \qquad
  \dd m' = \frac{\dd m}{M_0},
  \qquad
  \dd r'_i = \frac{r_{k+1} - r_k}{R_0},
\end{equation}
where $\tilde{S}_0 = S_0$ is a normalization constant (so $S' = S$ in
natural units of $\kbbar$).  $Y$ and $Z$ are already dimensionless
mass fractions and enter the state vector without rescaling.

The dimensionless transport coefficients used in the residual and
Jacobian are:
\begin{align}
  \varphi'  &= \frac{\varphi}{M_0\,T_0},
  \label{eq:phi_prime}\\
  \Lambda'  &= \frac{4\pi r'^2\,\Lambda\,R_0}{S_0\,M_0},
  \label{eq:Lambda_prime}\\
  \varphi'_{\rm semi}
            &= \frac{\varphi_{\rm semi}}{M_0\,T_0},
  \label{eq:semi_phi_prime}\\
  Q'        &= \frac{Q_{\rm cell}}{M_0\,T_0\,S_0},
  \label{eq:Q_prime}\\
  \left(\frac{\pp U}{\pp Y}\right)'
            &= \frac{1}{S_0}\frac{\pp U}{\pp Y}\bigg|_{S,\rho},
  \qquad
  \left(\frac{\pp U}{\pp Z}\right)'
            = \frac{1}{S_0}\frac{\pp U}{\pp Z}\bigg|_{S,\rho},
  \label{eq:dUdYZ_prime}\\
  \left(\frac{\pp T}{\pp Y}\right)'
            &= \frac{1}{T_0}\frac{\pp T}{\pp Y},
  \qquad
  \left(\frac{\pp T}{\pp Z}\right)'
            = \frac{1}{T_0}\frac{\pp T}{\pp Z}.
  \label{eq:dTdYZ_prime}
\end{align}
Here $\varphi$ is the convective transport prefactor
(Equations~\ref{eq:phi_env} or \ref{eq:phi_mantle}), $\Lambda$ is the
total thermal conductivity (Equation~\ref{eq:Frad}), $\varphi_{\rm semi}$
is the semiconvective transport coefficient
(Equation~\ref{eq:Fsemi}), and $Q_{\rm cell} = H_{\rm spec}\,\dd m$
is the per-cell radiogenic power.

\subsection{Spatial Discretization}\label{sec:grid}

The planet is divided into $N$ mass shells (cells), indexed
$k=0$ (outermost, surface) through $k=N-1$ (innermost, center).
Fluxes are evaluated at $N+1$ cell boundaries (subscript~$b$, indices
$0$ through~$N$); the $N-1$ interior boundaries (subscript~$i$,
indices $0$ through $N-2$) lie between adjacent cells.

Cell-centered quantities ($T$, $P$, $\rho$, $S$, $Y$, $Z$) are
mapped to interior boundaries via:
\begin{itemize}
  \item \textbf{Geometric mean} for $T$, $P$, $\rho$, $r$:
    \;$X_i = \sqrt{X_k\,X_{k+1}}$.
  \item \textbf{Arithmetic mean} for $Y$, $Z$, $S$:
    \;$X_i = (X_k + X_{k+1})/2$.
\end{itemize}
Radial gradients are approximated by centered finite differences:
\begin{equation}\label{eq:fd_grad}
  \frac{\dd X}{\dd r}\bigg|_i
  \approx \frac{X_{k+1} - X_k}{\Delta r_i},
  \qquad
  \Delta r_i = r_{k+1} - r_k,
\end{equation}
where $r_k$ are cell-centered radii (geometric means of the bounding
radii).

Boundary masks ensure that the stencil respects domain edges: for
cell $k=0$ only the right-hand flux contributes, and for cell
$k=N-1$ only the left-hand flux contributes.  The outermost boundary
($b=0$) carries the surface radiative flux
(Section~\ref{sec:bc}); the innermost boundary ($b=N$) carries zero
flux (Section~\ref{sec:bc_bottom}).

\subsection{State Vector and Residual}\label{sec:residual}

The state vector has $3N$ components, interleaving entropy,
helium, and metallicity for each cell:
\begin{equation}\label{eq:state_vec}
  \mathbf{z} = (S'_0, Y_0, Z_0,\; S'_1, Y_1, Z_1,\; \ldots,\;
                S'_{N-1}, Y_{N-1}, Z_{N-1}).
\end{equation}
The residual vector $\mathbf{B}(\mathbf{z})$ is assembled from the
implicit (backward-Euler) discretization of
Equations~(\ref{eq:entropy_master})--(\ref{eq:Z_master}).
The root $\mathbf{B}=\mathbf{0}$ yields the updated state at
$t+\Delta t$.

\subsubsection{Entropy Residual}

For cell~$k$ (row $3k$ of $\mathbf{B}$):
\begin{equation}\label{eq:BS}
  B_{3k}
  = \underbrace{(S'_k - S'_{k,\rm old})}_{\text{time derivative}}
  + \underbrace{\frac{\Delta t}{\dd m'_k\,T'_k}\,\Delta F_k}_{\text{flux divergence}}
  + \underbrace{\frac{L}{T_0\,S_0}\,\frac{\Delta\chi_k}{T'_k}}_{\text{latent heat}}
  + \underbrace{\frac{(\pp U/\pp Y)'_k\,(Y_k - Y_{k,\rm old})
        + (\pp U/\pp Z)'_k\,(Z_k - Z_{k,\rm old})}{T'_k\,T_0}}_{\text{composition coupling}}
  - \underbrace{\frac{\Delta t}{\dd m'_k\,T'_k}\,Q'_k}_{\text{radiogenic}},
\end{equation}
where $\Delta\chi_k = \chi_k^{n+1} - \chi_k^n$ is the melt-fraction
change (Section~\ref{sec:latent}).  The \textbf{net flux divergence}
across cell~$k$ is $\Delta F_k = F_{k+1/2} - F_{k-1/2}$, where each
interior-boundary flux is the sum of three components:
\begin{align}
  F_{{\rm conv},i}
    &= -\varphi'_i
       \left|\frac{\dd S'_{\rm eff}}{\dd r'}\right|_i^{n_{\rm conv}},
  \label{eq:Fconv_disc}\\
  F_{{\rm rad},i}
    &= \Lambda'_i\,
       \frac{T'_{{\rm trial},k+1} - T'_{{\rm trial},k}}
            {\dd r'_i},
  \label{eq:Frad_disc}\\
  F_{{\rm semi},i}
    &= -\varphi'_{{\rm semi},i}
       \left|\frac{\dd S'_{\rm semi}}{\dd r'}\right|_i.
  \label{eq:Fsemi_disc}
\end{align}
Here $\varphi'_i$, $\Lambda'_i$, and $\varphi'_{{\rm semi},i}$ are
the dimensionless transport coefficients
(Equations~\ref{eq:phi_prime}--\ref{eq:semi_phi_prime}) evaluated at
interior boundary~$i$.  The radiative flux uses the \emph{trial}
temperature (Section~\ref{sec:trial_temp}) rather than the old
hydrostatic temperature.  The sign convention in
Equation~(\ref{eq:Fconv_disc}) follows Equation~(\ref{eq:Fconv});
the absolute-value--power construction is smoothed by the soft-hinge
function described in Section~\ref{sec:mlt}.  The radiogenic
heating term $Q'_k$ is nonzero only in mantle cells
($k_{\rm core}\le k < k_{\rm core,Fe}$).

\subsubsection{Helium Residual}

For cell~$k$ (row $3k+1$ of $\mathbf{B}$):
\begin{equation}\label{eq:BY}
  B_{3k+1}
  = (Y_k - Y_{k,\rm old})
  + \frac{\Delta t}{\dd m_k}\Bigl[
      \underbrace{\sigma_{1,+}\,(Y_{k+1}-Y_k)
                - \sigma_{1,-}\,(Y_k-Y_{k-1})}_{\text{diffusion}}
    + \underbrace{\sigma_{2,+}\,\mathcal{A}_{Y,i+}
                - \sigma_{2,-}\,\mathcal{A}_{Y,i-}}_{\text{advection}}
  \Bigr],
\end{equation}
where $\sigma_{1,\pm}$ and $\sigma_{2,\pm}$ are the diffusive and
advective flux coefficients (Equations~\ref{eq:sigma1}--\ref{eq:sigma2})
evaluated at the right ($+$) and left ($-$) interior boundaries of
cell~$k$, respectively.  The advective coefficient includes the smooth
sign: $\sigma_{2,\pm} = \sigma_{2,i}\,\widetilde{\mathrm{sgn}}(\mathcal{A}_Y)$.
The activation products $\mathcal{A}_{Y,i\pm}$ are evaluated at the
corresponding boundaries from Equation~(\ref{eq:Yflux_new}).
Terms referencing $Y_{k-1}$ vanish for $k=0$, and terms referencing
$Y_{k+1}$ vanish for $k=N-1$.  Note that the composition residuals use
the \emph{physical} (not dimensionless) mass element $\dd m_k$.

\subsubsection{Metallicity Residual}

The $Z$ residual (row $3k+2$) has the same structure as the $Y$
residual:
\begin{equation}\label{eq:BZ}
  B_{3k+2}
  = (Z_k - Z_{k,\rm old})
  + \frac{\Delta t}{\dd m_k}\Bigl[
      \sigma_{1,Z,+}\,(Z_{k+1}-Z_k)
    - \sigma_{1,Z,-}\,(Z_k-Z_{k-1})
    + \sigma_{2,Z,+}\,\mathcal{A}_{Z,i+}
    - \sigma_{2,Z,-}\,\mathcal{A}_{Z,i-}
  \Bigr],
\end{equation}
with the metal diffusive and advective coefficients $\sigma_{1,Z}$,
$\sigma_{2,Z}$ and the dual-miscibility activation product
$\mathcal{A}_Z$ (Section~\ref{sec:metal_rain}).

\subsection{Jacobian Structure}\label{sec:jacobian}

The Jacobian $\mathbf{J} = \pp\mathbf{B}/\pp\mathbf{z}$ is a
$3N\times 3N$ matrix with a block-tridiagonal structure.  Each
$3\times 3$ block couples the $(S', Y, Z)$ triplet of a cell to
itself and its two neighbors, giving a total half-bandwidth of~5
(columns $3k-3$ through $3k+5$ for row $3k$).  This sparsity is
exploited by storing $\mathbf{J}$ in compressed sparse column (CSC)
format.

\subsubsection{Entropy Row ($3k$)}\label{sec:jac_entropy}

The diagonal entry is
\begin{equation}\label{eq:J_SS_diag}
  J_{3k,\,3k}
  = 1 + J_{{\rm lat},S,k}
  + \frac{\Delta t}{\dd m'_k\,T'_k}
    \sum_{i\in\{k-\frac{1}{2},\,k+\frac{1}{2}\}}
    \left[
      -\frac{\varphi'_i}{\dd r'_i}\,
       \frac{\pp f}{\pp S'}\bigg|_i
      -\frac{\Lambda'_i}{\dd r'_i}\,
       \frac{T'_k}{\Cv[,k]}\,S_0
      -\frac{\varphi'_{{\rm semi},i}}{\dd r'_i}
    \right],
\end{equation}
where $J_{{\rm lat},S,k}$ is the latent-heat Jacobian contribution
(Section~\ref{sec:latent}), $\pp f/\pp S'$ denotes the derivative
of the convective flux function (the soft-hinge bracket raised to
$n_{\rm conv}$) with respect to the local entropy, and the sum runs
over the two boundaries flanking cell~$k$.  The radiative term
carries the factor $T'_k/\Cv[,k]$ because $\pp T'/\pp S' = (T'/\Cv)\,S_0$.

The off-diagonal entropy--entropy coupling to the left neighbor is
\begin{equation}\label{eq:J_SS_km1}
  J_{3k,\,3(k-1)}
  = \frac{\Delta t}{\dd m'_k\,T'_k}
    \left[
      \frac{\varphi'_i}{\dd r'_i}\,
       \frac{\pp f}{\pp S'}\bigg|_i
    + \frac{\Lambda'_i}{\dd r'_i}\,
       \frac{T'_{k-1}}{\Cv[,k-1]}\,S_0
    + \frac{\varphi'_{{\rm semi},i}}{\dd r'_i}
    \right],
  \qquad (k>0),
\end{equation}
where all coefficients are evaluated at boundary $i=k-\tfrac{1}{2}$.
The right-neighbor entry $J_{3k,\,3(k+1)}$ has the same form with
$k-1\to k+1$ and boundary $i=k+\tfrac{1}{2}$.

The entropy--composition coupling to $Y$ (column $3k+1$) is
\begin{equation}\label{eq:J_SY}
  J_{3k,\,3k+1}
  = \frac{\Delta t}{\dd m'_k\,T'_k}
    \sum_i\left[
      \frac{\varphi'_i}{\dd r'_i}\,
       \left(\frac{\pp S}{\pp Y}\right)'_{\!\!i}
       \frac{\pp f}{\pp S'}\bigg|_i
    - \frac{\Lambda'_i}{\dd r'_i}\,
       \left(\frac{\pp T}{\pp Y}\right)'_{\!\!k}
    \right]
  + \frac{(\pp U/\pp Y)'_k}{T'_k\,T_0}
  + J_{{\rm lat},Y,k},
\end{equation}
where $(\pp S/\pp Y)'_i$ is the Ledoux entropy--helium derivative
at interior boundary~$i$.
The $Z$-coupling $J_{3k,\,3k+2}$ is analogous with $Y\to Z$.

The semiconvective flux contributes \emph{no} Jacobian terms (2026
July rework): with the Nusselt coefficients frozen at their
old-state, time-lagged values, $F_{\rm semi}$ is treated as a
lagged (Picard) term --- exact in the residual, absent from
$\mathbf{J}$ --- the same consistent inexact-Newton scheme used for
the convective-entrainment and composition-overshoot diffusivities.
The previous hand-derived entries were finite-difference--verified
to be correct at semiconvectively \emph{active} faces but carried
no gate on the $\max(\nabla_{\!S},0)$ bracket clamp, so at clamped
(marginally stable) faces the Jacobian asserted couplings the
residual did not possess, degrading Newton convergence precisely in
the flickering regime the scheme must survive.

When the implicit viscous transition is enabled in the mantle,
additional terms
$\pp F_{\rm conv}/\pp\chi\cdot\pp\chi/\pp(S',Y,Z)$ are computed
by finite-differencing the convective flux with respect to $\chi$ at
each mantle boundary, then propagating through
$\pp\chi/\pp T\cdot\pp T/\pp(S',Y,Z)$.

\subsubsection{Helium Row ($3k+1$)}\label{sec:jac_helium}

The diagonal entry is
\begin{equation}\label{eq:J_YY_diag}
  J_{3k+1,\,3k+1}
  = 1
  + \frac{\Delta t}{\dd m_k}\Bigl[
      -\sigma_{1,+} - \sigma_{1,-}
      + \text{(advection product-rule terms)}
    \Bigr],
\end{equation}
where the advection terms are combinations of
$\sigma_{2,\pm}\,\widetilde{\mathrm{sgn}}\,\mathcal{P}_1/2$
and $\sigma_{2,\pm}\,\widetilde{\mathrm{sgn}}\,\mathcal{P}_2/2$
arising from the product-rule expansion of
$\pp\mathcal{A}_Y/\pp Y$ (Equation~\ref{eq:dAY_dY}).

The helium--entropy coupling (column $3k$) encodes the
miscibility-curve temperature dependence:
\begin{equation}\label{eq:J_YS}
  J_{3k+1,\,3k}
  = \frac{\Delta t}{\dd m_k}\,\sigma_{2,i}\,
    \widetilde{\mathrm{sgn}}(\mathcal{A}_Y)\,\mathcal{P}_2\,
    \left(\frac{\pp\Yeq}{\pp T}\right)_{\!\!i}
    \frac{T_i}{\Cp[,i]}\,\frac{S_0}{2},
\end{equation}
with opposite signs for the left and right boundary contributions.
Off-diagonal Y--Y couplings to neighbors $k\pm1$ contain the
diffusive $\pm\sigma_1$ terms and advective corrections from
$\pp\mathcal{A}_Y/\pp Y_{k\pm1}$.

\subsubsection{Metallicity Row ($3k+2$)}\label{sec:jac_metal}

The $Z$ Jacobian has the same structure as the $Y$ Jacobian but
involves derivatives of both miscibility curves.  The
entropy--metallicity coupling includes terms proportional to
$\pp Z_{1,{\rm misc}}/\pp T$ and $\pp Z_{2,{\rm misc}}/\pp T$,
each multiplied by the corresponding $\softplus$ factor from the
dual activation product $\mathcal{A}_Z$
(Section~\ref{sec:metal_rain}).

\subsubsection{Surface Jacobian Terms}

For the outermost cell ($k=0$), the surface radiative-loss Jacobian
entries (Equations~\ref{eq:Jsurf_S}--\ref{eq:Jsurf_Z}) are added to
the entropy row.  These couple the surface entropy, helium, and
metallicity to the atmosphere boundary condition
(Section~\ref{sec:bc}).

\subsection{Newton--Raphson Iteration}\label{sec:newton}

The coupled $3N\times 3N$ system $\mathbf{B}(\mathbf{z})=\mathbf{0}$ is
solved by a damped Newton iteration.  At each iterate~$n$:
\begin{equation}\label{eq:NR}
  \mathbf{J}^{(n)}\,\delta\mathbf{z}
  = \mathbf{B}^{(n)},
  \qquad
  \mathbf{z}^{(n+1)}
  = \mathbf{z}^{(n)} - w_{\rm eff}\,\delta\mathbf{z},
\end{equation}
where $w_{\rm eff}\le 1$ is the effective damping weight after
backtracking (Section~\ref{sec:trust}).  The sparse linear system is
assembled as a CSC matrix and solved via
\texttt{scipy.sparse.linalg.spsolve} (a direct sparse LU
factorization).

The full algorithm proceeds as follows:
\begin{enumerate}
  \item \textbf{Initialize.}
    Set $\mathbf{z}^{(0)} = (S'_{0,\rm old}, Y_{0,\rm old},
    Z_{0,\rm old}, \ldots)$ from the previous timestep.
    Set the damping weight $w = w_{\rm config}$ (typically 1),
    convergence error $\epsilon=1$, iteration counter $n=0$,
    and restart counter $r=0$.

  \item \textbf{Iterate} while
    $\epsilon > \epsilon_{\rm tol}$ (and
    $\epsilon_{\rm flux} > \epsilon_{\rm flux,tol}$ if
    high-precision mode is active):
    \begin{enumerate}
      \item[(a)] Compute the trial temperature
        $T'_{\rm trial}$ from the current iterate
        (Section~\ref{sec:trial_temp}).
      \item[(b)] Evaluate all transport coefficients:
        convective prefactor~$\varphi'$,
        thermal conductivity~$\Lambda'$,
        semiconvective coefficient~$\varphi'_{\rm semi}$,
        diffusion coefficients $D_Y$, $D_Z$,
        advective coefficients $D_2$, $D_{2Z}$,
        miscibility curves $\Yeq$, $Z_{1,2,\rm misc}$,
        activation products $\mathcal{A}_Y$, $\mathcal{A}_Z$,
        and melt fraction~$\chi$.
      \item[(c)] Assemble $\mathbf{J}$ and $\mathbf{B}$ as described in
        Sections~\ref{sec:residual}--\ref{sec:jacobian}.
      \item[(d)] Convert $\mathbf{J}$ to sparse CSC format and solve
        $\mathbf{J}\,\delta\mathbf{z} = \mathbf{B}$.
        If $\delta\mathbf{z}$ contains non-finite entries, raise a
        failure and proceed to step~3.
      \item[(e)] Apply the damped step with trust-region backtracking
        (Section~\ref{sec:trust}) to obtain $\mathbf{z}^{(n+1)}$
        and the effective step $w_{\rm eff}\,\delta\mathbf{z}$.
      \item[(f)] Evaluate the convergence error
        (Section~\ref{sec:convergence}).
      \item[(g)] Update: set $\mathbf{z}^{(n)} \leftarrow \mathbf{z}^{(n+1)}$;
        extract $S' = z_{::3}$, $Y = z_{1::3}$, $Z = z_{2::3}$.
        Clip any negative $Z$ entries to zero.
        Increment~$n$; break if $n > 150$.
    \end{enumerate}

  \item \textbf{On failure} (non-finite $\delta\mathbf{z}$, or
    backtracking exhausted):
    increment $r$.  If $r > 3$, raise a fatal error.
    Otherwise, halve the damping weight ($w \leftarrow w/2$) and
    restart from $\mathbf{z}^{(0)}$.

  \item \textbf{Post-process.}
    Convert back to physical units:
    $S_{\rm new} = S'\,S_0/\tilde{S}_0$.
    Multiply all flux arrays by $M_0\,T_0\,S_0$ to recover
    physical units (erg\,s$^{-1}$).
    Optionally enforce core entropy monotonicity
    ($S_{k,\rm new} \le S_{k,\rm old}$ for $k \ge k_{\rm core,Fe}$).
\end{enumerate}

\subsubsection{Trust-Region Guards}\label{sec:trust}

Given the Newton correction $\delta\mathbf{z}$, the step
$\mathbf{z}_{\rm cand} = \mathbf{z}^{(n)} - w\,\delta\mathbf{z}$
is tested against three guards.  Starting from $w = w_{\rm config}$
and halving up to 10 times:
\begin{enumerate}
  \item \textbf{Finiteness:} all entries of $\mathbf{z}_{\rm cand}$
    must be finite (no NaN or $\pm\infty$).
  \item \textbf{Entropy floor:} all entropy entries must satisfy
    $S'_k > S_{\rm floor}$, where $S_{\rm floor} = 10^{-4}$.
  \item \textbf{Relative entropy cap:} the per-cell relative change
    must satisfy
    \begin{equation}\label{eq:trust_cap}
      \frac{|\Delta S'_k|}{\max(|S'_k|,\;S_{\rm floor})}
      < 0.25
      \qquad \forall\;k.
    \end{equation}
    This limit prevents violent melt-fraction and viscosity swings
    near the solidification front, where $\chi(T)$ is steep.
\end{enumerate}
If all three guards pass, the candidate is accepted with
$w_{\rm eff} = w$.  If 10 halvings are exhausted without
satisfying the guards, a failure is raised and the solver
proceeds to the restart logic (step~3 above).

\subsubsection{Convergence}\label{sec:convergence}

Iteration terminates when the relative step norm falls below
a tolerance:
\begin{equation}\label{eq:conv_crit}
  \epsilon
  = \max_i \frac{|w_{\rm eff}\,\delta z_i|}
                 {\max(|z_i|,\;z_{\rm floor})}
  < \epsilon_{\rm tol},
\end{equation}
with $z_{\rm floor} = 10^{-10}$ and a typical
$\epsilon_{\rm tol}\sim10^{-10}$.  In high-precision mode, an
additional flux-convergence criterion is enforced:
\begin{equation}\label{eq:flux_conv}
  \epsilon_{\rm flux}
  = \max_j\frac{|F_j^{(n)} - F_j^{(n-1)}|}
               {\max(|F_j^{(n)}|,\;|F_j^{(n-1)}|)}
  < \epsilon_{\rm flux,tol},
\end{equation}
where the maximum runs over all flux arrays (convective, radiative,
semiconvective, latent, helium diffusion/advection, and metal
diffusion/advection).

\subsection{Trial Temperature}\label{sec:trial_temp}

At each Newton iterate, a trial temperature is constructed by
linearizing around the previous hydrostatic state:
\begin{equation}\label{eq:Ttrial}
  T'_{\rm trial}
  = T' + \frac{T'}{C_{V,{\rm safe}}}\,S_0\,(S' - S'_{\rm old})
  + \left(\frac{\pp T'}{\pp Y}\right)(Y - Y_{\rm old})
  + \left(\frac{\pp T'}{\pp Z}\right)(Z - Z_{\rm old}),
\end{equation}
where $T'$ and the partial derivatives $\pp T'/\pp Y$,
$\pp T'/\pp Z$ are evaluated from the last accepted hydrostatic
state, and $C_{V,{\rm safe}} = \max(C_V,\,10^{-99})$ guards
against division by zero in degenerate zones.

The trial temperature is \emph{not} simply a diagnostic; it
enters the solver in three critical ways:
\begin{enumerate}
  \item \textbf{Melt fraction.}
    $\chi(T_{\rm trial})$ is recomputed from
    Equation~(\ref{eq:chi}) and relaxed toward the previous $\chi$
    (Equation~\ref{eq:chi_relax}).  This determines the latent-heat
    source term and the viscous transport geometry in the mantle.
  \item \textbf{Radiative flux.}
    The temperature difference
    $\Delta T'_{\rm trial} = T'_{{\rm trial},k+1} - T'_{{\rm trial},k}$
    enters the radiative flux
    (Equation~\ref{eq:Frad_disc}), making the conductive heat
    transport respond to the evolving state within each Newton step.
  \item \textbf{Miscibility activation.}
    The implicit temperature correction $\delta_{\rm imp}$ in the
    activation product (Equation~\ref{eq:Yflux_new}) depends on
    $S - S_{\rm old}$, which is driven by the trial temperature.
\end{enumerate}
This coupling ensures that the Newton solver sees a self-consistent
thermodynamic state at each iteration, rather than lagging behind
the hydrostatic update.

% ======================================================================
\section{Eddy Overshoot}\label{sec:overshoot}
% ======================================================================

An optional artificial overshoot transport augments the thermal
conductivity at envelope boundaries near the envelope--mantle interface:
\begin{equation}
  \kappa_{\rm ov}(r)
  = f_{\rm ov}\,v_{\rm MLT,anchor}\,H_{P,{\rm anchor}}
    \,\exp\!\left(-\frac{|r - r_{\rm EMB}|}{f_{H_P}\,H_{P,{\rm EMB}}}\right),
\end{equation}
where $v_{\rm MLT,anchor}$ is the convective velocity at the
strongest-bracket mantle boundary and $f_{\rm ov}$ is a user-specified
efficiency factor.  The resulting conductive enhancement
$\Lambda_{\rm ov} = \rho\,\Cp\,\kappa_{\rm ov}$ is capped at a
configurable multiple of the local conductivity.

% ======================================================================

\begin{acknowledgments}
ORCHARD builds on the APPLE code \citep{Sur2024}.
\end{acknowledgments}

\bibliography{transport_methods}
\bibliographystyle{aasjournal}

\end{document}
